\documentclass[conference]{IEEEtran}

\usepackage{amsmath,amssymb,amsthm}
\usepackage{booktabs}
\usepackage{multirow}
\usepackage{array}
\usepackage{enumitem}
\usepackage{hyperref}
\usepackage{graphicx}
\usepackage{float}
\usepackage{dblfloatfix}
\usepackage{placeins}
\usepackage{xcolor}
\usepackage{algorithm}
\usepackage{algpseudocode}
\usepackage{caption}
\usepackage{xspace}
\usepackage{balance}
\usepackage[most]{tcolorbox}

\newtcolorbox{rqbox}{
  colback=gray!5,
  colframe=black!70,
  boxrule=0.5pt,
  arc=2pt,
  left=4pt,
  right=4pt,
  top=2pt,
  bottom=2pt
}

\hypersetup{colorlinks=true,linkcolor=blue,citecolor=blue,urlcolor=blue}
\captionsetup{font=small}
\setlist[itemize]{leftmargin=1.5em}
\setlist[enumerate]{leftmargin=1.8em}
\setlength{\textfloatsep}{8pt plus 1pt minus 2pt}
\setlength{\dbltextfloatsep}{8pt plus 1pt minus 2pt}
\setlength{\floatsep}{6pt plus 1pt minus 2pt}
\setlength{\intextsep}{6pt plus 1pt minus 2pt}

\newcommand{\tool}[1]{\texttt{#1}}
\newcommand{\code}[1]{\texttt{#1}}
\newcommand{\iacpg}{IACPG\xspace}

\title{IACPG: Interrupt-Aware Code Property Graph for Agent-Based Defect Detection in Embedded Interrupt-Driven Programs}
\author{\IEEEauthorblockN{Author}
\IEEEauthorblockA{\textit{Affiliation}\\
email@example.com}}

\begin{document}
\maketitle

\begin{abstract}
While Large Language Model agents show great promise for automated defect detection, applying them to safety-critical embedded programs requires structured context beyond raw source code. Although the Code Property Graph (CPG) is widely employed to fuse syntax and data dependencies into a unified representation, it remains fundamentally inadequate for interrupt-driven scenarios. Specifically, the CPG fails to capture essential semantics such as asynchronous preemption, priority nesting, and masking constraints, leaving the agent without the explicit evidence required to reason about complex concurrency interactions. To address this, we propose the Interrupt-Aware Code Property Graph (IACPG). Designed specifically to empower the reasoning process of the agent, IACPG overlays interrupt roles, preemption relationships, guard constraints, and cross-context shared access edges onto the CPG. Evaluated on 96 cross-platform instances spanning four defect patterns and four instruction set architectures (ISAs), the agent powered by IACPG achieves 95.6\% precision, 97.7\% recall, and 96.6\% F1-score, outperforming the CPG-based agent by 8.0 percentage points and a zero-shot LLM by 6.8 percentage points in F1-score. Compared to the CPG, IACPG reduces context length and end-to-end response time by approximately 88.7\% and 88.4\%, respectively, while improving evidence completeness from 6.2\% to 86.5\%.
\end{abstract}

\begin{IEEEkeywords}
Defect detection, Agent-based analysis, Code Property Graph, Interrupt semantics, Embedded systems 
\end{IEEEkeywords}

\section{Introduction}
Interrupt-driven programs are widely used in embedded scenarios such as industrial control, automotive electronics, aerospace, and smart terminals. Their reliability and security directly impact the normal execution of physical devices and critical tasks \cite{martinez2014hardware, pan2019easy}. In such systems, the interrupt mechanism is extensively employed to handle peripheral events, meet real-time response requirements, and coordinate system execution flows, forming a core interrupt-driven execution paradigm where the main program and Interrupt Service Routines (ISRs) run cooperatively \cite{brylow2004deadline}. Unlike traditional sequential programs, the behavior of interrupt-driven programs is not entirely determined by explicit control flow. The asynchronous triggering of ISRs, priority nesting, and the scope of interrupt enabling and masking jointly shape the true execution order and shared state access semantics of the program.

These semantics pose fundamental challenges to code analysis. Interrupt concurrency defects (e.g., race conditions, atomicity violations) are often concealed and difficult to reproduce \cite{li2023empirical}. Once triggered during actual execution, they can lead to severe consequences such as data inconsistency, system crashes, or loss of control over physical devices. Although existing research has made significant progress in program graph modeling \cite{yamaguchi2014modeling, mi2023graph, liu2023learning, shezan2023chkplug}, defect detection \cite{eslamimehr2018efficient, zhang2025effectively, qiu2024vulnerability, smaili2025transformer, ma2025value, wang2025sift}, and embedded concurrency analysis \cite{wang2017automatic, li2022precise, yu2023detecting, wang2022specchecker, wu2016static, zhang2025bounded, zhao2025formal, wu2015numerical, schwarz2011static, shi2015verifying, yu2018conpredictor, wang2020automatic, liang2017effective, ma2026integrating}, it still struggles to meet the analysis needs in interrupt-driven scenarios. On the one hand, mainstream program representations like AST, CFG, PDG, and CPG are primarily designed for sequential programs or threaded concurrency scenarios and lack the direct ability to express interrupt-driven semantics \cite{zhu2023code}. Furthermore, many embedded concurrency analysis methods rely on specific architectural knowledge, manual rules, or dedicated detection logic, making it difficult to balance transferability, unified modeling capabilities, and result interpretability \cite{feng2020rchecker, wang2025specchecker}.

In recent years, Large Language Model (LLM)-driven code agents have shown great potential in program comprehension, defect localization, automated auditing, and code review \cite{li2024enhancing, liu2025llm, wang2024llmdfa, nguyen2025empirical, li2025codedoctor, yu2024fine, nashaat2024towards}. However, the quality of such analysis heavily relies on whether the underlying program representation can provide sufficient and structured semantic support \cite{lekssays2025llmxcpg, shao2025graphfvd, peng2025keep}. An intuitive approach is to supply raw source code directly to the agent to detect interrupt concurrency defects. However, raw source code lacks explicit semantic and structural information, leading to suboptimal detection effectiveness. To mitigate this structural deficiency, introducing the Code Property Graph (CPG) serves as a logical progression. As a well-established representation in traditional static program analysis, it fuses syntax, control flow, and data dependencies into a unified graph structure. By querying the CPG, the agent can extract a cleaner and more structured context compared to raw source code. Nevertheless, the CPG remains fundamentally inadequate for interrupt-driven scenarios: it cannot capture essential interrupt semantics such as asynchronous preemption, priority nesting, and masking constraints. Consequently, even with CPG-based context, the agent still lacks the targeted, structured evidence required to accurately reason about complex concurrency interactions. The inadequate expressive power of existing program representations for interrupt-driven semantics makes it difficult for agents to accurately recover potential interactions between interrupt contexts. Even if suspicious access patterns are found, they lack a unified structure to link the execution context, semantic constraints, and defect evidence, thereby affecting analysis accuracy and result interpretability. In summary, interrupt concurrency defect detection faces the following challenges: First, interrupt execution semantics (e.g., ISR identity, preemption relationships, interrupt masking scope, cross-context shared access) lack explicit modeling and are difficult to reflect directly in a unified representation. Second, existing methods struggle to provide structured, interpretable evidence chains to support the reasoning and decision-making processes of the agent while maintaining transferability.

Based on the above observations, this paper argues that the key to enabling agent-based analysis of interrupt-driven programs lies in constructing a unified program representation that can explicitly express interrupt execution semantics. While retaining the basic structure of the CPG, this representation should explicitly incorporate key semantic facts such as ISR identity, interrupt switches, priority levels, and shared variable accesses. It should further construct interrupt roles, preemption relationships, switch constraints, and cross-context shared access relationships, forming an enhanced program representation suitable for agent querying, reasoning, and interpretation. Based on this representation, the analysis system can more specifically recover potential interactions between execution contexts, narrow down the candidate space for interrupt concurrency issues, and form a more complete structured evidence chain. To this end, this paper proposes IACPG, an interrupt-semantics-enhanced code property graph method for the agent-based analysis of embedded interrupt-driven programs.


The main contributions of this paper are as follows:
\begin{itemize}
  \item \textbf{A Hybrid Interrupt Semantic Extraction Framework:} We propose a robust extraction method that combines a deterministic rule-based knowledge base with structurally constrained LLM completion. This approach enables the accurate recovery of key interrupt semantics (e.g., ISR identity, priorities, masking scopes) for the target architectures evaluated in this work, with restricted completion for uncovered callback-style cases.
  
  \item \textbf{The Interrupt-Aware Code Property Graph (IACPG):} We formulate a non-destructive semantic overlay on the CPG that explicitly encodes interrupt roles, preemption relationships, guard constraints, and cross-context shared accesses. By providing an agent-friendly structured representation, IACPG fundamentally streamlines the agent reasoning process and significantly reduces query load.
  
  \item \textbf{Controlled Cross-Platform Evaluation:} We evaluate the proposed approach on 96 instances spanning four defect patterns and four ISAs. The results demonstrate that IACPG not only substantially reduces reasoning hops and context length but also improves end-to-end defect detection accuracy, false-positive control, and diagnostic evidence completeness compared to existing baselines.
\end{itemize}

\section{Background and Problem Definition}

\subsection{Execution Characteristics of Interrupt-Driven Programs}

Interrupt-driven concurrency arises from hardware events, vector binding, priority scheduling, and interrupt masking rather than explicit thread creation: any main-program segment can be preempted by an asynchronously triggered ISR, and higher-priority handlers may nest within lower-priority ones. Without atomicity protection or masking constraints, shared variables are vulnerable to inconsistent interleaving, so effective analysis must explicitly recover execution context identity, preemption relationships, local interrupt switch constraints, and cross-context shared-state accesses.

\subsection{Deficiencies of Traditional Program Representations}

The Code Property Graph (CPG) unifies AST, CFG, and PDG into a single representation and has become the foundation for many program analysis tools. However, CPG is designed for sequential or thread-parallel execution models and lacks direct support for interrupt-driven semantics: it cannot describe which contexts an ISR might preempt, whether a program point is interruptible, or whether a shared access truly involves cross-context interference. This gap has two consequences: (1)~retrieving shared-variable accesses via conventional data flow yields large volumes of irrelevant candidates, expanding the search space; (2)~even when true concurrency defects are detected, the absence of a unified evidence chain linking execution context, interrupt constraints, shared access, and defect judgment limits result interpretability--a particular bottleneck for AI agents relying on structured context.

\subsection{Problem Definition}

Given an embedded C project $P$ and its code property graph $G_{cpg}$, we seek to construct an interrupt-aware enhanced graph $G_{iacpg}$ that explicitly models shared-state contention, preemption nesting, and masking constraints:
\begin{equation}
G_{iacpg} = \Phi(G_{cpg}, \mathcal{R}_{int})
\end{equation}
where $\mathcal{R}_{int}$ denotes the interrupt semantic enhancement items derived from the interrupt semantic facts $\mathcal{F}$ extracted from $P$, and $\Phi$ injects interrupt roles, preemption relationships, guard constraints, and cross-context shared access relationships via a semantic overlay while preserving the base CPG topology.




\section{Methodology}

\begin{figure*}[!htbp]
  \centering
  \includegraphics[width=0.92\linewidth]{Fig1.pdf}
  \caption{IACPG overall workflow and module organization.}
  \label{fig:overall-architecture-placeholder}
\end{figure*}

The proposed framework (Figure~\ref{fig:overall-architecture-placeholder}) follows a three-stage pipeline: \textbf{Stage~I} extracts interrupt semantic facts from source code using a hybrid rule-based plus LLM-completion approach; \textbf{Stage~II} injects these facts into the CPG to produce \iacpg; and \textbf{Stage~III} deploys an LLM agent that queries \iacpg via structured graph query interfaces and a feature-aware verification routing strategy to detect interrupt concurrency defects with traceable evidence chains.

\subsection{Interrupt Semantic Extraction}

Interrupt semantic extraction recovers, from embedded C source, four categories of facts forming the semantic foundation of \iacpg: execution context identities, interrupt control operations, ISR priority levels, and cross-context shared variable accesses. Because naming conventions, driver frameworks, and platform interfaces vary widely, pure static rules lack coverage while free LLM inference lacks verifiability. We therefore use a hybrid architecture: an engineering rule knowledge base as the primary path, with structurally constrained LLM completion as a controlled supplement; all outputs must pass CPG structural legality verification before entering the fact set $\mathcal{F}$.

\subsubsection{Rule Knowledge Base-Driven Extraction}

The rule knowledge base encodes platform-level ISR naming conventions, vector registration formats, interrupt enable/mask interfaces, priority configuration methods, and shared variable access patterns for known embedded platforms, focusing on architecture-level conventions and common implementation patterns. It serves as the primary extraction path, providing deterministic, high-confidence facts. Only when rule matching yields no result or encounters a conflict is the restricted agent completion path triggered.

\subsubsection{Restricted Completion and Structural Legality Verification}

For scenarios beyond rule coverage (e.g., deep platform encapsulation, indirect registration, complex macro expansions), an LLM agent completes missing facts. Concretely, the rule coverage check $\Call{RuleCoverage}{\mathcal{F}_{base}}$ in Algorithm~\ref{alg:build-iacpg} is treated as \emph{insufficient} when at least one ISR-candidate function---identified by the isolated in-degree pattern (no incoming call edges from ordinary functions) in $G_{cpg}$---carries no rule-matched fact in any of the three semantic dimensions ISR/Mask/Priority. The restricted completion path is then triggered on those uncovered functions only, rather than on the whole project, keeping the agent invocation scope bounded. Agent output is strictly constrained to four semantic categories (ISR, main context, interrupt control, shared access) and must pass two structural legality checks before merging: (1)~\textbf{Node Existence}: classified targets must exist in the CPG node set; (2)~\textbf{Topology Constraint}: ISR candidates must satisfy the isolated in-degree pattern (no incoming call edges from ordinary functions), and shared access candidates must be anchored to actual read/write nodes appearing in multiple contexts. Spurious ISRs not satisfying these constraints are suppressed during the downstream relationship derivation stage.

\subsection{IACPG Graph Construction}
  
Figure~\ref{fig2} depicts the transformation from source code to IACPG. On top of the base CPG, the recovered semantic facts are injected as attribute enrichment ($\mathcal{R}_{role}$ tags on function nodes) and three new edge types: GUARD (interrupt-masking boundaries), PREEMPT (cross-context execution interferences), and SHARED ACCESS (conflicting access points of shared variables). This externalizes implicit hardware-level constraints as explicit graph structures while preserving base-CPG compatibility.

Formally, let the embedded project be $P$, and the CPG be
\begin{equation}
G_{cpg}=(V,E,\tau,\lambda)
\end{equation}
where $V$ represents the set of nodes, $E$ represents the set of edges, $\tau$ represents the node type mapping, and $\lambda$ represents the node attribute mapping. We recover a semantic fact set related to interrupt-driven execution on top of $G_{cpg}$:
\begin{equation}
\mathcal{F}=\mathcal{F}_{ctx}\cup \mathcal{F}_{ctrl}\cup \mathcal{F}_{prio}\cup \mathcal{F}_{acc}
\end{equation}
where $\mathcal{F}_{ctx}$, $\mathcal{F}_{ctrl}$, $\mathcal{F}_{prio}$, and $\mathcal{F}_{acc}$ represent execution context facts, interrupt control facts, priority facts, and shared access facts, respectively. The priority hierarchy undergoes architecture-aware normalization to eliminate the heterogeneity of underlying hardware priority definitions, ensuring the logical layer uniformly follows the semantic constraint that higher numerical values correspond to higher preemption weights.

Based on these facts, a set of interrupt semantic enhancement items $\mathcal{R}_{int}$ is derived to characterize interaction constraints between execution flows:
\begin{equation}
\mathcal{R}_{int} = \mathcal{R}_{role}\cup \mathcal{R}_{preempt}\cup  \mathcal{R}_{guard}\cup \mathcal{R}_{share}
\end{equation}
where $\mathcal{R}_{role}$ assigns execution context role tags to nodes; $\mathcal{R}_{preempt}$, $\mathcal{R}_{guard}$, and $\mathcal{R}_{share}$ describe cross-context preemption possibilities, local switch constraints, and potential shared access conflicts, respectively.

Thus, \iacpg is defined as the result of the CPG mapped through the semantic enhancement mapping $\Phi$:
\begin{equation}
G_{iacpg} = \Phi(G_{cpg}, \mathcal{R}_{int}) = (V, E \cup E_{int}, \tau, \lambda')
\end{equation}

This non-destructive construction establishes a Semantic Overlay Network anchored to the base graph. Candidate interactions are represented as quadruples $q=(c_s,c_t,v,\Gamma)$, where $c_s$, $c_t$ are execution contexts, $v$ is the shared variable, and $\Gamma$ captures priority and guard conditions. A candidate enters structured analysis only when $q$ satisfies: distinguishable roles, reachable preemption path, unguarded by local constraints, and cross-context shared access.

\subsubsection{Interrupt Roles, Preemption, and Guard Constraints}

Function nodes are tagged with execution roles (main context, ISR, ordinary function). For main-context/ISR interactions, a \textit{worst-case concurrency assumption} conservatively treats any reachable, non-globally-masked ISR as capable of preempting any main-context point, ensuring recall completeness. For ISR nesting, \textit{platform-aware deduction} checks architectural rules in $\mathcal{K}$: only when the target platform supports nesting and the normalized priority satisfies $j > i$ is an ISR-nests-ISR edge added, preventing cross-platform false positives. For guard constraints ($\mathcal{R}_{guard}$), intra-procedural CFG reachability approximates masking scopes; ambiguous cross-function protection regions are delegated to the downstream agent for inter-procedural verification.

\subsubsection{Shared Access Modeling and Construction Process}

Cross-context shared access edges are added when two contexts both access variable $v$ and a preemption path exists between them. Read/write attributes are preserved to support fine-grained race and atomicity analysis. The full construction pipeline---\textbf{Fact Recovery} $\to$ \textbf{Relationship Derivation} $\to$ \textbf{Graph Injection} $\to$ \textbf{Analysis Readiness}---is formalized in Algorithm~\ref{alg:build-iacpg}. \textsc{IsUnsafeInteraction} applies a triple filter (role heterogeneity, path reachability, guard coverage) at the access-pair level, connecting $\mathcal{R}_{share}$ edges only when an unprotected preemption possibility exists---enabling detection of complex temporal patterns such as atomicity violations ($n_1 \to n_3 \gets n_2$).

\begin{algorithm}[!b]
\caption{Algorithm for Building IACPG Based on Interrupt Semantic Recovery}
\label{alg:build-iacpg}
\begin{algorithmic}[1]
\Require Embedded C project source $P$, CPG extractor $\mathcal{B}$, architecture \& rule knowledge base $\mathcal{K}$, Agent $\mathcal{A}$
\Ensure Semantic overlay enhanced property graph $G_{iacpg}$

\State $G_{cpg} \gets \mathcal{B}(P)$
\State $\mathcal{F}_{base} \gets \Call{RuleExtract}{P, G_{cpg}, \mathcal{K}}$

\If{$\Call{RuleCoverage}{\mathcal{F}_{base}}$ is insufficient}
  \State $\Delta\mathcal{F}_{raw} \gets \mathcal{A}.\Call{InferSemantics}{P, G_{cpg}, \mathcal{F}_{base}}$
  \State $\Delta\mathcal{F}_{valid} \gets \Call{StructValidate}{\Delta\mathcal{F}_{raw}, G_{cpg}}$
  \State $\mathcal{F} \gets \mathcal{F}_{base} \cup \Delta\mathcal{F}_{valid}$
\Else
  \State $\mathcal{F} \gets \mathcal{F}_{base}$
\EndIf

\State $\mathcal{R}_{role} \gets \Call{MapRoles}{\mathcal{F}}$
\State $\mathcal{R}_{preempt} \gets \Call{PlatformAwarePreempt}{\mathcal{F}, \mathcal{K}}$
\State $\mathcal{R}_{guard} \gets \Call{IntraProcGuards}{\mathcal{F}, G_{cpg}}$
\State $\mathcal{R}_{share} \gets \emptyset$

\For{each variable $v \in \mathcal{F}_{acc}$}
  \For{each pair of access nodes $(n_i, n_j)$ on $v$} 
    \If{\Call{IsUnsafeInteraction}{$n_i, n_j, \mathcal{R}_{preempt} \cup \mathcal{R}_{guard}$}}
      \State $\mathcal{R}_{share} \gets \mathcal{R}_{share} \cup \{(n_i, n_j, v)\}$
    \EndIf
  \EndFor
\EndFor

\State $\mathcal{R}_{int} \gets \mathcal{R}_{role} \cup \mathcal{R}_{preempt} \cup \mathcal{R}_{guard} \cup \mathcal{R}_{share}$
\State $G_{iacpg} \gets \Call{ApplyOverlayNetwork}{G_{cpg}, \mathcal{R}_{int}}$

\State \Return $G_{iacpg}$
\end{algorithmic}
\end{algorithm}

\begin{figure*}[!htbp]
  \centering
  \includegraphics[width=0.9\linewidth]{Fig2.pdf}
  \caption{IACPG Graph Construction}
  \label{fig2}
\end{figure*}

\subsection{Agent-Based Defect Detection via IACPG}

In this work, the agent serves as an LLM-based diagnostic verifier that issues structured graph queries and validates candidate interactions over IACPG, rather than a general-purpose autonomous software engineering agent. The LLM does not freely explore the graph; instead, it operates over route-selected primitive outputs and synthesizes the final evidence chain from these structured retrieval results.
The agent queries $G_{iacpg}$ through four \textbf{Semantic Retrieval Primitives}: (1)~\textit{Contextual Topology Discovery} retrieves execution roles from $\mathcal{R}_{role}$; (2)~\textit{Preemption Path Mapping} extracts cross-context jump relationships from $\mathcal{R}_{preempt}$; (3)~\textit{Guard Domain Delimitation} locates masking scopes via $\mathcal{R}_{guard}$ reachability; (4)~\textit{Interaction-aware Slicing} retrieves read/write sequences and data dependencies for a specific shared variable. This transforms agent reasoning from passively ingesting source code to proactively issuing targeted graph queries, reducing context load while ensuring evidence traceability. In RQ2, each reasoning hop corresponds to one such structured retrieval step (i.e., one primitive invocation) rather than one free-form conversational turn.

A \textbf{feature-aware verification routing strategy} then dispatches candidate interactions $q=(c_s,c_t,v,\Gamma)$ to specialized channels: \textit{Composite Timing State Conflict} analysis (for Read-Modify-Write patterns and multi-word accesses) verifies that $c_t$ can preempt the critical window via $\mathcal{R}_{preempt}$, has write capability on $v$, and that no $\mathcal{R}_{guard}$ protection covers the window; \textit{Data Dependency Pollution} analysis (for index/divisor boundary defects) verifies that the validation-to-use path is preemption-exposed and that $c_t$ can invalidate the boundary condition. Routing is deterministic at the defect-family level: atomicity-violation and multi-word-data-race candidates trigger the former channel, whereas interrupt-aware array OOB and divide-by-zero candidates trigger the latter. For example, when validating an atomicity-violation candidate, the agent first retrieves role topology, then checks preemption reachability, next queries guard coverage, and finally inspects a shared-access slice before synthesizing the final evidence package. Upon confirmation, the system synthesizes a traceable evidence package comprising: execution role topology, minimal violation trace, and protection gap identification---providing structured input for automated repair or manual review.






\FloatBarrier
\section{Experiment Design}

\subsection{Research Questions}

We answer three research questions:
\begin{itemize}
  \item \textbf{RQ1:} Can the proposed hybrid parsing method accurately recover key semantic facts such as interrupt entry points, priorities, masking constraints, and shared accesses in heterogeneous platform code?
  \item \textbf{RQ2:} Compared to CPG, does the explicit modeling of interrupt semantics in \iacpg significantly reduce the query load and reasoning hops for agents during defect mining?
  \item \textbf{RQ3:} Compared to baseline methods without interrupt semantic enhancement, can \iacpg effectively improve end-to-end defect detection performance while providing complete, interpretable evidence chains?
\end{itemize}

\subsection{Dataset Construction and Evaluation Targets}

\begin{table}[!t]
  \centering
  \caption{Benchmark statistics by defect type (96 instances, 4 ISAs $\times$ 24 templates).}
  \label{tab:dataset_src}
  \resizebox{\columnwidth}{!}{%
  \begin{tabular}{lcccc}
    \toprule
    \textbf{Defect Type} & \textbf{Inst.} & \textbf{Defective} & \textbf{\#Bugs} & \textbf{\#ISR} \\
    \midrule
    Atomicity Violation (AV)  & 24 & 24 & 24 & 32 \\
    Array Out-of-Bounds (OOB) & 24 & 16 & 16 & 24 \\
    Divide-by-Zero (DBZ)      & 24 & 24 & 24 & 32 \\
    Multi-word Data Race (DR) & 24 & 24 & 24 & 28 \\
    \midrule
    \textbf{Total}                  & \textbf{96} & \textbf{88} & \textbf{88} & \textbf{116} \\
    \bottomrule
  \end{tabular}}
  \\[4pt]
  \footnotesize{Note: simple\_001--004 are public benchmark seeds; simple\_005--006 are manually curated, architecture-adapted, and LLM-assisted extensions derived from real-world code patterns.}
\end{table}

The evaluation set (Table~\ref{tab:dataset_src}) comprises 96 instances across four ISAs (ARM, RISC-V, AVR, MSP430) and four defect patterns; seed sources are described below.

\begin{itemize}
  \item \textbf{Atomicity Violation (AV):} Templates (simple\_001--004) derived from the interrupt-driven RaceBench 2.1 benchmark~\cite{racebench2021repo}, also used in prior atomicity-violation studies such as intAtom~\cite{li2022precise}, containing 31 C programs with hand-injected atomicity violations.
  \item \textbf{Interrupt-aware Array OOB and Divide-by-Zero:} Templates (simple\_001--004) adapted from the cases released with SpecChecker-Int~\cite{wang2025specchecker}, covering concurrency-induced runtime errors in interrupt-driven programs.
  \item \textbf{Multi-word Data Race (DR):} Templates (simple\_001--004) constructed from real-world embedded code fragments exhibiting multi-word shared-variable races.
  \item \textbf{Extended instances (simple\_005--006 per defect type):} Manually curated, architecture-adapted, and LLM-assisted extensions derived from real-world code patterns to broaden architectural coverage.
\end{itemize}

The 8 non-defective instances are BufferOverflow templates free of interrupt-concurrent defects and serve as negative controls; each instance carries 1--2 ISRs and at most one shared variable (0.92 on average).


\begin{table*}[!t]
  \centering
  \caption{RQ1: Accuracy analysis of interrupt semantic fact extraction (based on 96 program instances).}
  \label{tab:rq1}
  \renewcommand{\arraystretch}{1.1}
  \resizebox{\textwidth}{!}{
  \begin{tabular}{l *{6}{cccccc} c}
    \toprule
    & \multicolumn{6}{c}{\textbf{Static Analysis (Ours)}} & \multicolumn{6}{c}{\textbf{LLM Baseline}} & \\
    \cmidrule(lr){2-7}\cmidrule(lr){8-13}
    \textbf{Semantic Dimension} 
    & \textbf{TP} & \textbf{FP} & \textbf{FN} 
    & \textbf{Precision} & \textbf{Recall} & \textbf{F1} 
    & \textbf{TP} & \textbf{FP} & \textbf{FN} 
    & \textbf{Precision} & \textbf{Recall} & \textbf{F1} 
    & \textbf{$\Delta$F1} \\
    \midrule
    ISR Identification 
    & 112 & 0 & 4 & 1.000 & 0.965 & 0.982 
    & 95  & 1 & 21& 0.990 & 0.819 & 0.896 
    & $+$0.086 \\
    
    IRQ Mask           
    & 118 & 0 & 0 & 1.000 & 1.000 & 1.000 
    & 98  & 24& 20& 0.803 & 0.831 & 0.817 
    & $+$0.183 \\
    
    Priority           
    & 25  & 0 & 2 & 1.000 & 0.926 & 0.962 
    & 18  & 0 & 9 & 1.000 & 0.667 & 0.800 
    & $+$0.162 \\
    
    Shared Access      
    & 88  & 0 & 0 & 1.000 & 1.000 & 1.000 
    & 75  & 0 & 13& 1.000 & 0.852 & 0.920 
    & $+$0.080 \\
    \midrule
    \textbf{Average}   
    &   &   &   & \textbf{1.000} & \textbf{0.973} & \textbf{0.986} 
    &   &   &   & \textbf{0.948} & \textbf{0.792} & \textbf{0.858} 
    & \textbf{$+$0.128} \\
    \bottomrule
  \end{tabular}
  }
  \begin{quote}
    \footnotesize
    Note: ISR = Interrupt Service Routine; IRQ Mask = Interrupt Request Mask; $\Delta$F1 = Static Analysis F1 $-$ LLM Baseline F1. The LLM baseline directly extracts the four semantics without tools or rule constraints.
  \end{quote}
\end{table*}

\subsection{Baselines and Evaluation Metrics}

Three configurations are compared: \textbf{LLM Zero-shot} (direct LLM query without graph support), \textbf{Agent (CPG)} (agent querying Joern CPG), and \textbf{Agent + IACPG} (our method). We additionally report SpecChecker-Int as the external specialized baseline in RQ3 because it is publicly available and covers the interrupt-induced runtime-error families most closely aligned with our benchmark, whereas tools such as intAtom, NIChecker, and CPA4AV focus on narrower defect scopes or do not provide directly compatible public artifacts for our four-category end-to-end setting. This comparison is intended to contrast defect detection effectiveness on the shared benchmark rather than equal-cost automation, since SpecChecker-Int is a fully automatic specialized static tool while IACPG uses LLM-assisted structured verification. Metrics are organized by RQ: (RQ1)~Precision and F1-score for four semantic extraction tasks; (RQ2)~Context Length, End-to-End Time, Reasoning Hops, Evidence Completeness (EC), and Signal-to-Noise Ratio (SNR); (RQ3)~Precision, Recall, and F1-score per ISA and overall (micro-averaged), with label-set agreement, defined as the ratio of instances whose predicted defect-label set exactly matches the ground-truth label set, as a supplementary overprediction metric. EC and SNR are used here as representation-oriented diagnostic metrics: they measure whether the graph structure surfaces the evidence elements required for interrupt-concurrency defect detection, rather than serving as universal benchmark-quality measures. EC in particular should be read as a representation-oriented coverage metric tied to the three interrupt-specific evidence elements, rather than a universal evidence-quality measure, and we accordingly avoid using it as a standalone headline claim. EC checks whether a report covers ISR identification, shared variable access, and concurrency window; SNR is the ratio of interrupt-relevant evidence characters to total tool-call output. All LLM-based configurations use the same MiniMax-M2.7 backend through an Anthropic-compatible API, with no task-specific decoding override applied in the evaluation scripts; thus, group differences primarily reflect representation and workflow differences rather than model substitution.

\section{Evaluation}

\subsection{RQ1: Accuracy of Interrupt Semantic Extraction}

Table~\ref{tab:rq1} compares our static analysis method against a direct LLM baseline across 96 instances. Under known-architecture rule base coverage, the deterministic extractor remains near-perfect across all four semantic dimensions, with F1 scores of 0.982 for ISR Identification, 1.000 for IRQ Mask, 0.962 for Priority, and 1.000 for Shared Access. The LLM baseline degrades significantly: ISR Identification F1=89.6\% (FP=1, FN=21, caused by missing architectural constraints), IRQ Mask F1=81.7\% (FP=24 from hallucinated switch operations), and Priority F1=80.0\% (FN=9 from platform-specific API failures). Average static-analysis F1 is 0.986, and average $\Delta$F1 remains +0.128. Importantly, these results are reported after removing benchmark-specific callback-name heuristics from the static rule base. The remaining misses are concentrated on a small set of non-standard callback-style ISR wrappers, which we intentionally treat as out-of-coverage for deterministic static rules rather than absorbing through case-specific name matching. These cases motivate the restricted agent completion path, but they should not be interpreted as a direct measurement of out-of-coverage generalization in the present benchmark.

To probe the fallback branch more directly, we additionally performed a controlled rule-holdout stress test. For each ISA, we removed only the ISR-identification patterns from the rule base while keeping switch and priority rules unchanged, and compared three settings on the same ISA subset: full rule coverage (Group~A), held-out static extraction only (Group~B), and held-out extraction with restricted completion (Group~C). Table~\ref{tab:rq1_holdout} shows that the completion path substantially restores ISR recognition under held-out conventions: ISR F1 recovers from 0.000 to 1.000 on ARM, from 0.585 to 1.000 on MSP430, from 0.000 to 1.000 on AVR, and from 0.000 to 0.951 on RISC-V. Priority recovery also returns from 0.000 to 1.000 on ARM and RISC-V. Shared Access remains at 1.000 across all groups because the current shared-variable extractor primarily depends on cross-function global-access evidence rather than strict ISR role labeling, making it a non-sensitive dimension in this stress test. Cases such as \texttt{on\_timer\_event} (MSP430) and \texttt{timer\_tick\_handler} (AVR) show that the completion path can recover genuinely non-standard callback- or macro-encoded ISR forms under controlled out-of-coverage conditions. This holdout test validates fallback utility under controlled rule removal within our benchmark platforms, but does not substitute for genuine unseen-platform evaluation.

\begin{table}[t]
  \centering
  \caption{RQ1 supplementary controlled rule-holdout stress test on ISR identification. FRB: full rule base; B: held-out static only; HORC: held-out + restricted completion.}
  \label{tab:rq1_holdout}
  \resizebox{\columnwidth}{!}{
  \begin{tabular}{lcccc}
    \toprule
    \textbf{ISA} & \textbf{FRB F1} & \textbf{HOSO F1} & \textbf{HORC F1} & \textbf{Recovery} \\
    \midrule
    ARM    & 0.982 & 0.000 & 1.000 & 1.018 \\
    AVR    & 0.982 & 0.000 & 1.000 & 1.018 \\
    MSP430 & 0.982 & 0.585 & 1.000 & 1.044 \\
    RISC-V & 0.982 & 0.000 & 0.951 & 0.968 \\
    \bottomrule
  \end{tabular}
  }
  \begin{quote}
    \footnotesize
    Note: Recovery is computed as $(F1_C - F1_B) / (F1_A - F1_B)$. Values above 1 indicate that Group~C slightly exceeds the current deterministic baseline on the same ISA subset. Priority recovery is additionally observed on ARM (0.963 $\rightarrow$ 0.000 $\rightarrow$ 1.000) and RISC-V (0.960 $\rightarrow$ 0.000 $\rightarrow$ 1.000).
  \end{quote}
\end{table}

\begin{rqbox}
\textbf{Answer to RQ1:}
Static extraction remains highly accurate under configured rule coverage (avg.\ F1 = 0.986), clearly exceeding the direct LLM baseline (avg.\ F1 = 0.858). The controlled rule-holdout stress test further shows that restricted completion effectively restores held-out ISR conventions across all four ISAs, confirming the fallback branch's utility for out-of-coverage patterns.
\end{rqbox}

% \noindent\textbf{Answer to RQ1:} Static extraction remains highly accurate under configured rule coverage (avg.\ F1 = 0.986), clearly exceeding the direct LLM baseline (avg.\ F1 = 0.858). The controlled rule-holdout stress test further shows that restricted completion effectively restores held-out ISR conventions across all four ISAs, confirming the fallback branch's utility for out-of-coverage patterns.

\subsection{RQ2: Structured Representation and Query Efficacy Gain}

RQ2 compares Agent(CPG) and Agent(IACPG) on the same set of benchmark instances using the same MiniMax-M2.7 backend. It contrasts a Joern-CPG retrieval workflow with an IACPG-enabled retrieval workflow across two dimensions: \textit{query efficacy} (context length, reasoning hops) and \textit{evidence quality} (EC, SNR). We further clarify that this comparison evaluates the combined effect of the interrupt-aware representation and its corresponding structured retrieval workflow (the four semantic retrieval primitives and the feature-aware verification routing), rather than an isolated representation-only ablation; the reported gains should therefore be read as benefits of the IACPG-enabled retrieval workflow as a whole.

As summarized in Table~\ref{tab:rq2}, IACPG demonstrates significant advantages across all dimensions. In the revised time accounting, response time is reported as end-to-end wall-clock time: IACPG includes Stage~1/2/3 preprocessing, while CPG includes the initial \texttt{joern\_import}. For IACPG, this end-to-end time also subsumes the underlying CPG construction performed inside Stage~3 via \texttt{joern-parse}/\texttt{joern-export}; by contrast, context length, reasoning hops, EC, and SNR remain online retrieval/evidence metrics measured after the representation is available. Under this mixed but explicit measurement scheme, IACPG reduces context length by $-$88.7\%, end-to-end response time by $-$88.4\%, and reasoning hops by $-$82.7\%, while improving EC by $+$80.3\,pp and SNR by $+$17.0\,pp.

\begin{table}[H]
  \centering
  \caption{RQ2: CPG vs.\ \iacpg in query efficacy and evidence quality.}
  \label{tab:rq2}
  \resizebox{\columnwidth}{!}{
  \begin{tabular}{lccc}
    \toprule
    \textbf{Evaluation Metric} & \textbf{CPG} & \textbf{IACPG}& \textbf{Improvement} \\
    \midrule
    \multicolumn{4}{l}{\textit{Query Efficiency}} \\
    Context Length (Chars)   & 21,640 & 2,438 & $-$88.7\% \\
    End-to-End Time (ms)    & 33,930 & 3,935 & $-$88.4\% \\
    Avg. Reasoning Hops  & 29.8 & 5.2 & $-$82.7\% \\
    \midrule
    \multicolumn{4}{l}{\textit{Evidence Quality}} \\
    Evidence Completeness (EC \%)    & 6.2\% & 86.5\% & $+$80.3\,pp \\
    Signal-to-Noise Ratio (SNR \%)   & 2.56\% & 19.57\% & $+$17.0\,pp \\
    \bottomrule
  \end{tabular}
  }
  \begin{quote}
    \footnotesize Note: End-to-end time includes representation construction (IACPG Stage~1/2/3, including the underlying CPG build inside Stage~3; CPG \texttt{joern\_import}). Reasoning hops, SNR, and EC are online retrieval/evidence metrics.
  \end{quote}
\end{table}

\begin{rqbox}
\textbf{Answer to RQ2:}
IACPG substantially improves the combined representation-and-retrieval workflow: compared with CPG, it shortens context, reduces reasoning hops, and improves interrupt-evidence quality. These gains indicate that interrupt-aware semantic edges allow the agent to reach defect-relevant evidence with much less search noise.
\end{rqbox}

% \noindent\textbf{Answer to RQ2:} IACPG substantially improves the combined representation-and-retrieval workflow: compared with CPG, it shortens context, reduces reasoning hops, and improves interrupt-evidence quality. These gains indicate that interrupt-aware semantic edges allow the agent to reach defect-relevant evidence with much less search noise.

\subsection{RQ3: Cross-Platform End-to-End Defect Detection Performance}

For positive instances we check the primary defect label; for negative controls (\texttt{defect\_classes=[]}) all four labels are exhaustively checked (any detection = FP). As shown in Table~\ref{tab:end2end_ase}, Agent + IACPG achieves the highest overall F1 of 96.6\% (P=95.6\%, R=97.7\%), outperforming Agent (CPG) at 88.6\%, LLM Zero-shot at 89.8\%, and SpecChecker-Int at 76.1\%. The gap is most pronounced on MSP430 (IACPG 95.5\% vs.\ CPG 78.9\%, $\Delta$=+16.6pp), where masking and priority semantics are most platform-specific. SpecChecker-Int's lower recall (67.0\%) stems from its conservative masking-scope inference: when a global disable/enable pair is detected, the tool treats the entire task as protected and suppresses conflict reports, leading to systematic false negatives on MDR and complex AV cases.

Although Zero-shot slightly exceeds CPG in F1 (89.8\% vs.\ 88.6\%), a supplementary label-set agreement analysis reveals its most severe overprediction: its exact predicted defect-label set matches the ground-truth label set on only 47.9\% of instances (vs.\ 63.5\% for CPG, 75.0\% for IACPG), with 37/88 positive instances carrying off-target extra labels. IACPG mitigates both false positives and over-reporting by providing precise interrupt-aware graph abstractions. Remaining IACPG errors: 2 spurious AV predictions on BufferOverflow/simple\_004 (ARM, RISC-V) and 2 spurious OOB predictions on BufferOverflow/simple\_006 (AVR, MSP430) due to approximate interrupt-masking scope inference; 2 FN on MultiwordDataRace/simple\_006 for AVR and MSP430. These 4 false positives are distributed over 4 of the 8 negative-control instances, while the remaining 4 negative instances incur no false positives.

\begin{rqbox}
\textbf{Answer to RQ3:}
Agent + IACPG delivers the strongest end-to-end defect detection performance across all four ISAs, achieving 95.6\% precision, 97.7\% recall, and 96.6\% F1 overall. It improves over both Agent (CPG) and LLM Zero-shot while keeping false positives limited to 4 of the 8 negative-control instances and yielding the largest per-ISA gain on MSP430.
\end{rqbox}

\begin{table}[H]
\centering
\caption{Cross-platform end-to-end defect detection performance (96 instances, 4 ISAs).}
\label{tab:end2end_ase}
\renewcommand{\arraystretch}{1.1}
\resizebox{\columnwidth}{!}{
\begin{tabular}{llcccccc}
\toprule
\textbf{Architecture (ISA)} & \textbf{Configuration} & \textbf{TP} & \textbf{FP} & \textbf{FN} & \textbf{Precision} & \textbf{Recall} & \textbf{F1} \\
\midrule
\multirow{4}{*}{ARM}    & SpecChecker-Int & 15 & 2 & 7 & 0.882 & 0.682 & 0.769 \\
                        & LLM (Zero-shot) & 19 & 1 & 3 & 0.950 & 0.864 & 0.905 \\
                        & Agent (CPG)     & 21 & 1 & 1 & 0.955 & 0.955 & 0.955 \\
                        & \textbf{Agent + IACPG} & \textbf{22} & \textbf{1} & \textbf{0} & \textbf{0.957} & \textbf{1.000} & \textbf{0.978} \\
\addlinespace[0.5em]
\multirow{4}{*}{RISC-V} & SpecChecker-Int & 14 & 2 & 8 & 0.875 & 0.636 & 0.737 \\
                        & LLM (Zero-shot) & 21 & 1 & 1 & 0.955 & 0.955 & 0.955 \\
                        & Agent (CPG)     & 20 & 2 & 2 & 0.909 & 0.909 & 0.909 \\
                        & \textbf{Agent + IACPG} & \textbf{22} & \textbf{1} & \textbf{0} & \textbf{0.957} & \textbf{1.000} & \textbf{0.978} \\
\addlinespace[0.5em]
\multirow{4}{*}{AVR}    & SpecChecker-Int & 15 & 2 & 7 & 0.882 & 0.682 & 0.769 \\
                        & LLM (Zero-shot) & 18 & 1 & 4 & 0.947 & 0.818 & 0.878 \\
                        & Agent (CPG)     & 18 & 1 & 4 & 0.947 & 0.818 & 0.878 \\
                        & \textbf{Agent + IACPG} & \textbf{21} & \textbf{1} & \textbf{1} & \textbf{0.955} & \textbf{0.955} & \textbf{0.955} \\
\addlinespace[0.5em]
\multirow{4}{*}{MSP430} & SpecChecker-Int & 15 & 2 & 7 & 0.882 & 0.682 & 0.769 \\
                        & LLM (Zero-shot) & 17 & 1 & 5 & 0.944 & 0.773 & 0.850 \\
                        & Agent (CPG)     & 15 & 1 & 7 & 0.938 & 0.682 & 0.789 \\
                        & \textbf{Agent + IACPG} & \textbf{21} & \textbf{1} & \textbf{1} & \textbf{0.955} & \textbf{0.955} & \textbf{0.955} \\
\midrule
\midrule
\multirow{4}{*}{\textbf{Overall (Micro-Avg.)}} & SpecChecker-Int & 59 & 8  & 29 & 0.881 & 0.670 & 0.761 \\
                                               & LLM (Zero-shot) & 75 & 4  & 13 & 0.949 & 0.852 & 0.898 \\
                                               & Agent (CPG)     & 74 & 5  & 14 & 0.937 & 0.841 & 0.886 \\
                                               & \textbf{Agent + IACPG} & \textbf{86} & \textbf{4} & \textbf{2} & \textbf{0.956} & \textbf{0.977} & \textbf{0.966} \\
\bottomrule
\end{tabular}
}
  \begin{quote}
    {\footnotesize Note: SpecChecker-Int is evaluated family-wise because it does not natively distinguish atomicity violations from multi-word data races. Overall values are micro-averaged over aggregated TP/FP/FN across all 96 instances.}
  \end{quote}
\end{table}


\subsection{Ablation Study}

\begin{table}[!t]
  \centering
  \caption{Ablation results over extraction, representation, and verifier stages.}
  \label{tab:ablation}
  \setlength{\tabcolsep}{2.4pt}
  \resizebox{\linewidth}{!}{%
  \begin{tabular}{llccc}
    \toprule
    \textbf{Configuration} & \textbf{Metric} & \textbf{Full} & \textbf{Variant} & \textbf{Change} \\
    \midrule
    w/o Restricted Completion & Avg. extraction F1 & 1.000 & 0.986 & $-0.014$ \\
    w/o Rule-based Extraction & Avg. extraction F1 & 1.000 & 0.858 & $-0.142$ \\
    w/o LLM Verifier & End-to-end F1 & 0.966 & 0.726 & $-0.240$ \\
    w/o Interrupt-aware Overlay & End-to-end F1 & 0.966 & 0.886 & $-0.080$ \\
    w/o Structured Graph Evidence & End-to-end F1 & 0.966 & 0.898 & $-0.068$ \\
    \bottomrule
  \end{tabular}%
  }
  \\[4pt]
  \footnotesize{Note: ``w/o'' means ``without''. The full framework uses hybrid extraction and Agent + IACPG. The variants correspond to Pure Rule extraction, Pure LLM extraction, IACPG + Static Checker, Agent + CPG, and Zero-shot LLM, respectively.}
\end{table}

Table~\ref{tab:ablation} reorganizes the RQ1--RQ3 outputs into component-level ablations and adds a rule-only diagnosis ablation over IACPG. The goal is not to introduce another dataset, but to make explicit which stage is responsible for each observed gain. The \textit{w/o Restricted Completion} and \textit{w/o Rule-based Extraction} rows isolate the extraction stage, \textit{w/o LLM Verifier} tests whether IACPG evidence alone is sufficient without agent-side diagnostic reasoning, and the remaining two rows isolate the representation and graph-evidence interface available to the agent.

\textbf{Extraction layer.} The hybrid extractor combines deterministic platform rules with restricted completion. Removing only restricted completion yields a small drop in average extraction F1 (1.000 to 0.986), because rules already cover most target-architecture conventions and completion is triggered only for four non-standard callback-style ISR cases. An isolated end-to-end rerun of these four cases with hybrid Stage~1 preserves the primary MDR detection in all four cases, although one AVR case still emits an off-target AV label; thus, hybrid extraction stabilizes the primary label but does not eliminate all label-set overprediction. In contrast, replacing the hybrid extractor with a pure LLM extractor reduces average extraction F1 to 0.858, showing that unconstrained language-model inference is insufficient for stable recovery of masking and priority semantics. This validates the rule-first design: deterministic rules provide high-confidence coverage, while restricted completion is a bounded fallback rather than the primary extraction mechanism.

\textbf{Representation and evidence interface.} The diagnosis-stage ablations show that both graph structure and interrupt-specific evidence matter. Replacing the LLM verifier with a rule-only static checker over IACPG drops F1 from 0.966 to 0.726 (TP=57, FP=12, FN=31), showing that explicit interrupt evidence alone is not sufficient without a verifier capable of combining timing, guard, and defect-family constraints. Replacing IACPG with standard CPG reduces end-to-end F1 from 0.966 to 0.886, indicating that generic data-flow and call-graph evidence is not enough to support interrupt-concurrency reasoning. Removing structured graph evidence entirely and using zero-shot LLM reasoning yields 0.898 F1, which confirms that the model can often infer the primary defect family from source code, but lacks the explicit role, preemption, guard, and shared-access evidence needed for robust and interpretable diagnosis. The ablation therefore supports the central design choice of combining queryable interrupt semantics with a structured diagnostic verifier rather than relying on either static rules or unstructured source-code reasoning alone.

\FloatBarrier

\section{Threats to Validity}

% \subsection{Discussion and Threats to Validity}

\textbf{Why MSP430 shows the largest gain.} MSP430's enable/disable intrinsics (e.g., \texttt{\_\_enable\_interrupt}/\texttt{\_\_disable\_interrupt}) and priority handling diverge more from the assumptions implicit in CPG than ARM NVIC does, leaving CPG's generic data-flow edges least able to capture interrupt-specific protection scopes; IACPG's explicit guard and preemption relations directly remediate this gap.

\textbf{Benchmark coverage boundary.} The 96-instance benchmark is derived from known-architecture platforms largely covered by the configured rule base. In the current revision, we deliberately removed benchmark-specific callback-name heuristics that directly matched non-standard handlers such as \texttt{system\_tick\_callback}, \texttt{timer\_tick\_handler}, \texttt{on\_timer\_event}, and \texttt{plic\_irq\_dispatch}. As a result, RQ1 now reports near-perfect rather than perfect static extraction, with the remaining misses concentrated on these callback-style ISR wrappers. We treat them as motivating cases for the restricted completion path rather than reabsorbing them into the deterministic rule base through case-specific name matching. A controlled rule-holdout stress test further confirms that this completion path can substantially recover held-out ISR conventions under benchmark-derived out-of-coverage simulation, while removing the callback-name heuristics does not alter the end-to-end IACPG defect detection results reported in RQ3. Large-scale evaluation of the same completion mechanism on industrial codebases remains concrete future work.

\textbf{Benchmark limitations.} The current benchmark comprises 24 defect templates instantiated across 4 ISAs, which improves architectural coverage but does not amount to 96 fully independent codebases. Most instances are synthetic or semi-synthetic, only 8 cases serve as negative controls, and each instance contains at most one shared variable (0.92 on average). These properties make the benchmark suitable for isolating representation-level effects, but they also limit the direct external generalization of the reported numbers to large industrial firmware projects with multiple source files, richer aliasing structures, and denser interrupt interactions.

\textbf{Scope of the agent component.} The LLM component in the current system serves primarily as a structured diagnostic verifier: it selects retrieval primitives, validates candidate interactions, and synthesizes evidence chains over IACPG. The core contribution remains the interrupt-aware representation and its evidence-oriented verification workflow rather than autonomous agent capabilities.

\textbf{Known error sources.} The residual 4 FP cases originate from intra-procedural approximation of interrupt-masking scopes: cross-function protection regions are delegated to the downstream agent for inter-procedural verification, which can over- or under-approximate guard coverage. Improving cross-procedure guard-scope inference is a concrete direction for further reducing this error.

\FloatBarrier
\section{Related Work}

\subsection{Embedded Interrupt Concurrency Analysis}

Embedded interrupt concurrency analysis primarily focuses on race conditions, atomicity violations, and priority-related defects between main programs and ISRs. Empirical studies show that such defects are highly destructive in industrial embedded systems \cite{li2023empirical}. Prior work spans dynamic testing, static analysis, and BMC, including SDRacer \cite{wang2017automatic}, sequentialization-based analyses \cite{wu2015numerical, wu2016static}, Rchecker \cite{feng2020rchecker}, atomicity-oriented tools such as intAtom and NIChecker \cite{li2022precise, zhang2025bounded}, delayed-trigger / partial-order methods such as CPA4AV and DIDP \cite{yu2023detecting, zhao2025formal}, and related studies on interrupt usability, data sharing, and concurrency prediction \cite{pan2019easy, wang2022specchecker, wang2025specchecker, yu2018conpredictor}.

These methods often rely on specific hardware architectures (e.g., Z86 timing analysis \cite{brylow2004deadline}), dedicated abstract models, or strong manual rules, leaving room for cross-platform adaptability and unified expressive capability. Rather than targeting a specific defect type, this paper builds an interrupt-aware representation as unified support for rule validation, agent-based analysis, and structured derivation.

\subsection{Program Representation and Defect Analysis Based on Program Graphs}

Graph-based representations such as AST, CFG, and CPG have become central in modern program analysis. Since the introduction of CPG for vulnerability mining \cite{yamaguchi2014modeling}, graph representations have been applied to code attribute classification and readability analysis \cite{liu2023learning, mi2023graph}, fine-grained vulnerability detection and prediction \cite{zhu2023code, shao2025graphfvd, nadim2022leveraging}, and compliance or context-aware analysis \cite{shezan2023chkplug, lekssays2025llmxcpg}.

These methods mostly assume sequential or thread-parallel execution, which makes it difficult to depict embedded-specific mechanisms such as ISR priorities, Interrupt Mask Registers, and asynchronous preemption. Rather than redesigning a general program graph, we inject an interrupt semantic layer atop CPG to fill this gap.

\subsection{LLM and Agent-Enhanced Program Analysis}

With the evolution of Large Language Models (LLMs), researchers have increasingly explored LLM and agent techniques for program comprehension, review, and defect detection. Capability studies such as Empica \cite{nguyen2025empirical} report strong semantic reasoning potential but continuing difficulty with deep data dependencies. Tool-augmented analysis systems include LLift, LATTE, and LLMDFA \cite{li2024enhancing, liu2025llm, wang2024llmdfa}, while related work has also explored formal verification and automated review/refactoring support \cite{ma2026integrating, yu2024fine, li2025codedoctor, nashaat2024towards}.

However, these methods still face bottlenecks when handling low-level hardware interactions and concurrency logic, and without a stable structured intermediate representation they suffer from output uncertainty and poor repeatability on embedded hardware states and asynchronous interrupt mechanisms. Unlike prior work focusing on end-to-end automated analysis, this paper proposes a hybrid framework combining restricted rule validation with structured agent verification over an ``interrupt-aware'' program graph, providing more interpretable underlying semantic support for embedded defect analysis.


\section{Conclusion and Future Work}

We propose IACPG, an interrupt-semantics-enhanced code property graph designed for the agent-based defect detection of embedded interrupt-driven programs. By combining a restricted rule base with structured evidence verification, IACPG systematically recovers ISR entities, execution priority topologies, guard masking logic, and cross-context shared-state interactions, subsequently injecting these facts as a multi-dimensional semantic overlay atop the CPG. In an end-to-end evaluation across four instruction set architectures (ISAs) and 96 instances, the agent powered by IACPG achieves 95.6\% precision, 97.7\% recall, and 96.6\% F1-score. This performance outcompetes the CPG-based agent by 8.0 percentage points and a zero-shot LLM by 6.8 percentage points in F1-score. Furthermore, compared to the CPG, IACPG reduces context length by 88.7\% and end-to-end response time by 88.4\%, while concurrently improving evidence completeness from 6.2\% to 86.5\% and the signal-to-noise ratio (SNR) from 2.56\% to 19.57\%. Under the current deterministic rule base, static extraction remains near-perfect within the configured rule coverage, and a controlled rule-holdout stress test demonstrates that the restricted completion path can substantially recover held-out ISR conventions across the four target ISAs.

Future work will extend the rule base to multi-core heterogeneous processors and RTOS environments, integrate lightweight path-sensitive alias analysis to tighten false-positive bounds on interrupt protection scopes, and connect the structured semantic foundation of IACPG to SMT solvers for fully automated patch-level verification.



\bibliographystyle{IEEEtran}
\bibliography{ref}
\end{document}
